\documentclass[11pt,a4paper]{article}

\usepackage[T1]{fontenc}
\usepackage[utf8]{inputenc}
\usepackage{amsmath}
\usepackage{amssymb}
\usepackage{booktabs}
\usepackage{listings}
\usepackage{graphicx}
\usepackage[margin=2.5cm]{geometry}
\usepackage[htt]{hyphenat}
\usepackage[colorlinks=true,linkcolor=blue,urlcolor=blue]{hyperref}

\setlength{\emergencystretch}{3em}

\lstset{
  basicstyle=\ttfamily\small,
  columns=fullflexible,
  breaklines=true,
  frame=single,
  framesep=4pt,
  xleftmargin=6pt,
  showstringspaces=false,
  language=Python,
}

\newcommand{\op}[1]{\texttt{#1}}
\newcommand{\cd}[1]{c^{\dagger}_{#1}}
\newcommand{\cc}[1]{c^{\phantom{\dagger}}_{#1}}

\title{\textbf{TranCI}\\[0.4em]
\large Configuration-interaction calculations for a transition-metal $d$ shell\\[0.6em]
\normalsize User manual and physical reference}
\author{}
\date{}

\begin{document}

\maketitle

\begin{abstract}
\noindent
TranCI solves the many-body problem of a single transition-metal $d$ shell
exactly, by full configuration interaction (exact diagonalization) in the space
of $n_e$ electrons distributed over the ten $d$ spin-orbitals.  The Hamiltonian
combines the full Coulomb repulsion between the $d$ electrons, crystal fields of
arbitrary symmetry, spin--orbit coupling, and magnetic and exchange fields.
From the resulting many-body spectrum the code extracts level degeneracies,
excitation energies, expectation values of spin and orbital operators, orbital
occupations, correlation entropies, dynamical correlation functions, and
low-energy effective spin Hamiltonians.  This document describes the physics
that the code implements, the exact form of every operator it provides, and how
to drive it both from Python and from the graphical interface.
\end{abstract}

\tableofcontents
\newpage

%======================================================================
\section{Overview}
%======================================================================

\subsection{What the code solves}

A transition-metal ion in a solid, on a surface, or in a molecule carries an
open $d$ shell.  Because the $d$ orbitals are compact, the Coulomb repulsion
between the $d$ electrons is comparable to, or larger than, the crystal-field
splitting induced by the environment.  A single-particle (mean-field) picture is
then not sufficient: the multiplet structure --- Hund's rules, the ordering of
the $^3F$, $^3P$, $^1G$ \ldots terms, magnetic anisotropy generated by the
interplay of spin--orbit coupling and the crystal field --- is a genuinely
many-body effect.

TranCI solves the problem without approximation inside the $d$ shell.  It
constructs the full many-body Hamiltonian in the basis of all Slater
determinants with $n_e$ electrons in ten spin-orbitals, and diagonalizes it
exactly.  Because the largest space has only $\binom{10}{5}=252$ states, this is
cheap, and every eigenstate is available.

\subsection{Structure of a calculation}

Every calculation follows the same three steps.

\begin{enumerate}
\item \textbf{Load the operator library for a given occupation.}
      \op{get\_atom(ne)} returns an \op{Atom} object holding a dictionary
      \op{Atom.Operator} of \emph{many-body} matrices --- all of them already
      expressed in the $\binom{10}{n_e}$-dimensional basis of that electron
      number.
\item \textbf{Assemble the Hamiltonian yourself} as a linear combination of
      those matrices.  There is no fixed model: the Hamiltonian is whatever
      Hermitian combination you write down.
\item \textbf{Diagonalize and analyse}, via
      \op{Atom.get\_manifolds(H)}, which returns a \op{Lowest\_States} object
      exposing the spectrum, the degenerate manifolds and all the observables.
\end{enumerate}

\begin{lstlisting}
import sys; sys.path.append("<repo>/src")
from tranci.atom import get_atom

Atom = get_atom(ne=3)                   # d^3 ion, 120 many-body states
V  = Atom.Operator["Coulomb"]           # electron-electron repulsion
CF = Atom.Operator["z2"]                # uniaxial crystal field
LS = Atom.Operator["ls"]                # spin-orbit coupling
Sz = Atom.Operator["sz"]                # total S_z

H = 4*V + 0.3*CF + 0.05*LS + 0.01*Sz    # your Hamiltonian, in eV
M = Atom.get_manifolds(H)               # diagonalize

print(M.get_gs_degeneracy())            # (degeneracy, energy)
print(M.get_excitations())              # excitation energies [eV]
print(M.get_gs_projected_eigenvalues(Sz))
\end{lstlisting}

\noindent
\textbf{Units.}  All energies are in electronvolts (eV) everywhere in the
library.  The only exception is the file \op{SPECTRUM.OUT} written by the
graphical interface, which additionally reports energies in GHz and meV.

\subsection{Installation and imports}

The library is \emph{not} installed as a Python package.  Every script must
prepend the source directory to the module search path:

\begin{lstlisting}
import sys
sys.path.append("/path/to/tranci/src")
from tranci.atom import get_atom
\end{lstlisting}

The bundled examples locate the repository relative to the \emph{current working
directory}, not to the script file, so they must be run from inside their own
directory:

\begin{lstlisting}
cd examples/octahedral && python main.py
\end{lstlisting}

Running \op{python examples/octahedral/main.py} from the repository root fails on
import.  The notebooks in \op{notebooks/} resolve the repository the same way,
so their kernel must be started with the working directory in \op{notebooks/}.

Dependencies are \op{numpy}, \op{scipy} and \op{matplotlib} for the core
library; \op{PyQt5} and a \op{pdflatex} installation for the graphical
interface; and \op{jax} for the effective-Hamiltonian fitting.  \op{numba} is
used if present but is optional.

%======================================================================
\section{The model}
%======================================================================

\subsection{Single-particle space}
\label{sec:spbasis}

The $d$ shell has $\ell=2$, hence $2\ell+1=5$ orbital states labelled by the
magnetic quantum number $m=-2,-1,0,+1,+2$, each of which can hold a spin-up or a
spin-down electron.  The single-particle space therefore has ten spin-orbitals
$|m,\sigma\rangle$.  The code orders them with spin as the slow index and $m$ as
the fast index:

\begin{center}
\begin{tabular}{ccccccccccc}
\toprule
index      & 0 & 1 & 2 & 3 & 4 & 5 & 6 & 7 & 8 & 9\\
$m$        & $-2$ & $-1$ & $0$ & $+1$ & $+2$ & $-2$ & $-1$ & $0$ & $+1$ & $+2$\\
$\sigma$   & $\uparrow$ & $\uparrow$ & $\uparrow$ & $\uparrow$ & $\uparrow$
           & $\downarrow$ & $\downarrow$ & $\downarrow$ & $\downarrow$ & $\downarrow$\\
\bottomrule
\end{tabular}
\end{center}

\noindent
This ordering is recorded in the file \op{orbital.in} shipped with each operator
library, and it is the ordering in which every $10\times 10$ single-particle
matrix must be written when it is supplied by the user
(Sec.~\ref{sec:custom}).

\subsection{Many-body space}

For $n_e$ electrons the many-body basis consists of all occupation-number
states (Slater determinants)
\begin{equation}
  |n_0 n_1 \ldots n_9\rangle
  \;=\;\prod_{i:\,n_i=1}\cd{i}\,|0\rangle ,
  \qquad n_i\in\{0,1\},\qquad \sum_i n_i = n_e ,
\end{equation}
with the creation operators applied in increasing index order, so that fermionic
signs are fixed unambiguously.  The dimension is the binomial coefficient
$\binom{10}{n_e}$:

\begin{center}
\begin{tabular}{lccccccccc}
\toprule
$n_e$        & 1 & 2 & 3 & 4 & 5 & 6 & 7 & 8 & 9\\
$\binom{10}{n_e}$ & 10 & 45 & 120 & 210 & 252 & 210 & 120 & 45 & 10\\
\bottomrule
\end{tabular}
\end{center}

The occupation vectors, one per row, are stored in \op{basis.out} inside each
operator library, and are read into \op{Atom.basis} as a list of
\op{BaseMB} objects.  Occupations $n_e=1$ through $9$ are supported; $n_e=0$ and
$n_e=10$ are trivial closed-shell cases and are rejected.

\subsection{Where the matrices come from}

The many-body matrices are not built at run time.  They are precomputed once by
a C\texttt{++} generator (\op{src/tranci/cpplib/}) and stored, in sparse text
form, in \op{src/tranci/cilib/<ne>/}.  Each \op{.op} file starts with a header
line \op{\# SIZE = d} and then lists \op{i j real imag} for the non-zero
elements of a Hermitian matrix.  \op{get\_atom(ne)} simply reads that directory.
This is why loading an atom is instantaneous and why the physical content of the
Coulomb tensor (Sec.~\ref{sec:coulomb}) is fixed by the shipped data.

%======================================================================
\section{The Hamiltonian}
\label{sec:ham}
%======================================================================

The Hamiltonian is written by the user as a linear combination of the operators
in \op{Atom.Operator}.  A typical complete Hamiltonian reads
\begin{equation}
  \mathcal{H}
  \;=\; U\,\hat V_{ee}
  \;+\; \hat V_{\mathrm{CF}}
  \;+\; \lambda \sum_i \vec{l}_i\cdot\vec{s}_i
  \;+\; \mu_B\vec{B}\cdot(\vec{L}+2\vec{S})
  \;+\; \vec{J}\cdot\vec{S} ,
  \label{eq:H}
\end{equation}
All coefficients are in eV, with the single exception of the Coulomb prefactor
$U$, which is a dimensionless multiplier (Sec.~\ref{sec:coulomb}).  The rest of
this section gives the precise operator that each dictionary key corresponds
to.

\subsection{Electron--electron interaction}
\label{sec:coulomb}

The key \op{"Coulomb"} (equivalently \op{"vc"}) holds the full two-body
interaction inside the $d$ shell,
\begin{equation}
  \hat V_{ee}
  \;=\;\sum_{ijkl}\;\sum_{\sigma\sigma'}
  V_{ijkl}\;
  \cd{i\sigma}\cd{j\sigma'}\cc{k\sigma'}\cc{l\sigma} ,
  \label{eq:vee}
\end{equation}
with $i,j,k,l$ running over the five orbital indices and the spin sums running
over the two spin channels.  The interaction is spin-independent and rotationally
invariant, so it conserves $m_i+m_j=m_k+m_l$: only the $85$ symmetry-allowed
components are non-zero.  Note that Eq.~\eqref{eq:vee} carries no factor
$\tfrac12$; the tensor $V_{ijkl}$ read from file is correspondingly one half of
the conventional Coulomb matrix element.

The tensor is the standard Slater--Condon decomposition of the Coulomb kernel
for $\ell=2$,
\begin{equation}
  V_{ijkl}\;\propto\;\sum_{\kappa=0,2,4} F^{\kappa}\;
  c^{\kappa}(m_i,m_l)\,c^{\kappa}(m_j,m_k),
\end{equation}
where the $c^{\kappa}$ are the usual Gaunt coefficients for $\ell=2$ and the sum
runs over the even ranks $\kappa=0,2,4$ allowed by parity.  All $85$ tabulated
magnitudes are reproduced exactly by a single set of Slater integrals.  (The
sign convention of the stored table is tied to the
$\cd{}\cd{}\cc{}\cc{}$ operator ordering used by the generator and is not a
separate physical input.)

The physically meaningful statement is the one about the operator that the user
actually applies.  With a multiplier $u$, i.e.\ for
$u\cdot\op{Atom.Operator["Coulomb"]}$, the effective Slater and Racah parameters
are
\begin{equation}
  F^{0}=7.70\,u\ \mathrm{eV},\qquad
  F^{2}=4.06\,u\ \mathrm{eV},\qquad
  F^{4}=2.65\,u\ \mathrm{eV},
\end{equation}
\begin{equation}
  B=\frac{F^{2}}{49}-\frac{5F^{4}}{441}=52.9\,u\ \mathrm{meV},
  \qquad
  C=\frac{35F^{4}}{441}=210.3\,u\ \mathrm{meV},
  \qquad C/B = 3.98 ,
\end{equation}
in line with the textbook value $C/B\approx 4$ for $3d$ ions and with the atomic
ratio $F^{2}/F^{4}\simeq 1.53$.  These values are not inferred from the file but
measured directly from the computed $d^2$ multiplet spectrum, which they
reproduce to all printed digits (Sec.~\ref{sec:validation}).

\paragraph{Important: the prefactor is a multiplier, not $F^0$.}
The number you place in front of \op{Atom.Operator["Coulomb"]} multiplies this
\emph{whole fixed tensor}; it is dimensionless, not an energy.  Writing
\op{H = 4*V} therefore means
\begin{equation}
  F^{0}=30.8\ \mathrm{eV},\quad F^{2}=16.2\ \mathrm{eV},\quad F^{4}=10.6\ \mathrm{eV},
  \quad B=212\ \mathrm{meV},\quad C=841\ \mathrm{meV},
\end{equation}
which is a very strongly correlated regime: the $d^2$ term splittings are then of
order $3$~eV.  The same holds for the field labelled \op{U [multiplier]} in the
graphical interface: its default value $U=4$ scales the tensor by $4$.  Treat
$U$ as a dimensionless strength of the multiplet interaction and tune it so that
the term splittings match the system you model.  Free $3d$
ions have $B\approx 0.10$--$0.13$~eV, reduced by $20$--$40\%$ in a solid, which
corresponds to $u\approx 1.5$--$2.5$.

Note finally that $F^0$ only shifts the total energy at fixed $n_e$, so it never
affects a spectrum computed at fixed occupation --- only $B$ and $C$ matter.

\subsection{Crystal fields}
\label{sec:cf}

The environment of the ion --- ligands, a surface, a molecular cage --- lowers
the spherical symmetry of the free ion and splits the five $d$ orbitals.  TranCI
supplies a set of Cartesian crystal-field generators built from powers of the
single-particle orbital angular momentum:
\begin{equation}
  \op{x2}=\sum_i (l_x^2)_i,\quad
  \op{y2}=\sum_i (l_y^2)_i,\quad
  \op{z2}=\sum_i (l_z^2)_i,\quad
  \op{x4}=\sum_i (l_x^4)_i,\ \ \text{etc.}
\end{equation}
\begin{equation}
  \op{x2y2}=\sum_i \Big[ (l_xl_y)^2 + (l_yl_x)^2 \Big]_i .
\end{equation}

\begin{quote}
\textbf{These are sums of single-particle operators, not powers of the total
angular momentum.}  That is, $\op{x2}=\sum_i (l_x^2)_i \neq L_x^2$.  The
many-body squares $L^2$, $S^2$, $J^2$ are available separately as \op{"l2"},
\op{"s2"}, \op{"j2"} and are built as $L_x^2+L_y^2+L_z^2$ from the total
operators.  Confusing the two is the most common source of unexpected spectra.
\end{quote}

\paragraph{Uniaxial field ($\op{z2}$).}
$D\sum_i (l_z^2)_i$ has single-particle eigenvalues $m^2 D$, so it splits the
$d$ shell as
\begin{equation}
  E(d_{z^2})=0,\qquad
  E(d_{xz}),E(d_{yz})=D,\qquad
  E(d_{xy}),E(d_{x^2-y^2})=4D .
\end{equation}
For $D>0$ the $d_{z^2}$ orbital is stabilized.  This is the tetragonal /
axially-compressed geometry; $D<0$ reverses the ordering.

\paragraph{Rhombic field ($\op{x2}-\op{y2}$).}
$E\sum_i (l_x^2-l_y^2)_i$ breaks the remaining $xy$ degeneracy.  Its
single-particle spectrum is $\{\pm 2\sqrt{3},\,\pm 3,\,0\}\times E$.

\paragraph{Octahedral field ($\op{x4}+\op{y4}+\op{z4}$).}
The cubic invariant $O\sum_i(l_x^4+l_y^4+l_z^4)_i$ has exactly two distinct
single-particle eigenvalues: $18\,O$ (six-fold: the $t_{2g}$ triplet
$d_{xy},d_{xz},d_{yz}$, times spin) and $24\,O$ (four-fold: the $e_g$ doublet
$d_{z^2},d_{x^2-y^2}$, times spin).  For $O>0$ the $t_{2g}$ states lie lowest,
as in an octahedral ligand cage, and the ligand-field splitting is
\begin{equation}
  10\,Dq \;=\; 6\,O .
\end{equation}
For $O<0$ the ordering is inverted, corresponding to a tetrahedral or cubic
environment.

\paragraph{Trigonal field.}
The trigonal generator is the uniaxial field rotated so that its symmetry axis
points along the body diagonal $\hat n=(1,1,1)/\sqrt3$,
\begin{equation}
  \hat V_{\mathrm{tri}} = t\; R\,\Big[\textstyle\sum_i (l_z^2)_i\Big]\,R^{\dagger}
  = t \sum_i (\vec l_i\cdot\hat n)^2
  = \frac{t}{3}\sum_i \big(l_x+l_y+l_z\big)^2_i ,
  \qquad
  R=e^{\,i\phi L_z}e^{\,i\theta L_x},
\end{equation}
with $\cos\theta=1/\sqrt3$ and $\phi=\pi/4$, as implemented in
\op{hamiltonians.build\_hamiltonian}.  The identity with $(\vec l\cdot\hat n)^2$
has been verified numerically for the $(1,1,1)$ diagonal specifically.  Note the
factor $1/3$, which the interface label ``$((x+y+z)/\sqrt3)^2$'' reflects: the
operator is normalized so that its single-particle eigenvalues are again
$0,1,4$, exactly as for \op{z2}.  In your own scripts you can build any rotated
field the same way, using the total $L_\alpha$ matrices as generators.

\paragraph{Arbitrary crystal fields.}
Any one-body crystal field can be supplied directly as a $10\times 10$ matrix
and promoted to the many-body basis; see Sec.~\ref{sec:custom}.

\subsection{Spin--orbit coupling}
\label{sec:soc}

The key \op{"ls"} holds the one-body spin--orbit operator summed over electrons,
\begin{equation}
  \hat H_{\mathrm{SOC}} = \lambda \sum_i \vec{l}_i\cdot\vec{s}_i
  = \lambda \sum_i \Big[ (l_z s_z)_i + \tfrac{1}{2}(l_+ s_- + l_- s_+)_i \Big].
\end{equation}
For a single electron in the $d$ shell this operator has eigenvalues
$-\tfrac{3}{2}$ (four-fold, $j=\tfrac{3}{2}$, verified through
$\langle j^2\rangle = \tfrac{15}{4}$) and $+1$ (six-fold, $j=\tfrac{5}{2}$),
so a positive $\lambda$ places the $j=3/2$ multiplet lowest and the splitting is
\begin{equation}
  \Delta E_{5/2-3/2} = \tfrac{5}{2}\lambda .
\end{equation}
\paragraph{Use $\lambda>0$ for every filling.}
The input $\lambda$ is the \emph{one-electron} spin--orbit constant, which is
positive for all $d$ ions.  The familiar sign change of the \emph{effective}
many-body coupling $\lambda_{LS}=\pm\zeta/2S$ above half filling is not
something you put in: it emerges from the many-body problem, and the code
reproduces it automatically.  Diagonalizing $u\,\hat V_{ee}+\lambda\,\op{ls}$
with $\lambda>0$ gives, for the free-ion ground terms,
\begin{center}
\begin{tabular}{lcccc}
\toprule
$n_e$ & 2 & 3 & 7 & 8\\
term  & $^3F$ & $^4F$ & $^4F$ & $^3F$\\
$J$ (computed) & $2$ & $3/2$ & $9/2$ & $4$\\
Hund III & $|L{-}S|$ & $|L{-}S|$ & $L{+}S$ & $L{+}S$\\
\bottomrule
\end{tabular}
\end{center}
i.e.\ $J=|L-S|$ below half filling and $J=L+S$ above it, exactly as Hund's third
rule requires, with the same positive $\lambda$ throughout.  Typical magnitudes
are $\lambda\sim 10$--$100$~meV for $3d$ ions, growing rapidly down the periodic
table.

\subsection{Magnetic and exchange fields}

Two distinct field terms are available.

\paragraph{Zeeman field.}  Couples to the total magnetic moment,
\begin{equation}
  \hat H_{Z} = \vec{b}\cdot(\vec{L}+2\vec{S}),
\end{equation}
using the many-body $L_\alpha$ and $S_\alpha$ (keys \op{"lx"}\ldots,
\op{"sx"}\ldots).  The vector $\vec{b}$ is entered in eV, i.e.\ it already
includes the Bohr magneton: $b = \mu_B B$ with
$\mu_B=5.7884\times 10^{-5}\ \mathrm{eV/T}$, so $1$~T corresponds to
$b\simeq 5.79\times 10^{-5}$~eV.  The factor $2$ is the free-electron spin
$g$-factor.

\paragraph{Exchange field.}  Couples to the spin only,
\begin{equation}
  \hat H_{J} = \vec{J}\cdot\vec{S},
\end{equation}
which models the molecular field of a magnetically ordered substrate or of
neighbouring magnetic ions.  Because it does not act on the orbital sector, it
is the right tool for isolating spin physics from orbital physics.

\subsection{Hyperfine coupling}

A spin-$\tfrac12$ nucleus can be added with \op{hyperfine.add\_nucleus(atom)},
which doubles the Hilbert space to $|{\rm electronic}\rangle\otimes|I_z\rangle$
and provides $\vec{I}$ operators together with $\vec{S}\cdot\vec{I}$.  The
Hamiltonian then acquires
\begin{equation}
  \hat H_{\mathrm{hf}} = A\,\vec{S}\cdot\vec{I} + g_n\,\vec{b}\cdot\vec{I},
  \qquad g_n \simeq 1/1836 ,
\end{equation}
with the nuclear Zeeman term scaled by the proton-to-electron mass ratio.  This
path is exercised only through \op{build\_hamiltonian} when
\op{atom.has\_nucleus} is set; it is not exposed in the graphical interface.

%======================================================================
\section{Diagonalization, manifolds and degeneracies}
%======================================================================

\subsection{Exact diagonalization}

\op{Atom.get\_manifolds(H)} converts the sparse Hamiltonian to a dense matrix
and diagonalizes it completely with a Hermitian LAPACK routine.  All
$\binom{10}{n_e}$ eigenvalues and eigenvectors are computed and retained; no
truncation of the many-body space is made anywhere.  The cost scales as
$\dim^3$, which for the worst case $\dim=252$ is milliseconds.

The result is a \op{Lowest\_States} object with
\begin{itemize}
\item \op{.evals} --- the exact eigenvalues, \emph{shifted so that the ground
      state is at zero};
\item \op{.e0} --- the absolute ground-state energy before that shift;
\item \op{.evecs} --- the eigenvectors, one per row;
\item \op{.h} --- the Hamiltonian that was diagonalized;
\item \op{.atom} --- the originating \op{Atom} object.
\end{itemize}

\subsection{Degenerate manifolds}

Symmetry makes the spectrum highly degenerate, and almost every physical
statement (degeneracy of the ground state, magnitude of a zero-field splitting,
which excitations are allowed) requires grouping eigenvalues into manifolds.
Two eigenvalues are declared degenerate when
\begin{equation}
  |E_a - E_b| < \texttt{tol},
  \qquad \texttt{tol} = 10^{-\texttt{ntol}} ,
\end{equation}
where \op{tol} and \op{ntol} are \emph{module-level globals} of
\op{tranci.hamiltonians} (default $\op{ntol}=6$, i.e.\ $\op{tol}=10^{-6}$~eV).
The graphical interface overwrites them at run time from the
``Energy tolerancy'' field.  From a script you set them explicitly:

\begin{lstlisting}
from tranci import hamiltonians
hamiltonians.tol  = 1e-5       # grouping window in eV
hamiltonians.ntol = 5
\end{lstlisting}

\noindent
This window is a \emph{grouping} parameter, not a display precision:
\op{.evals} always holds the exact eigenvalues.  Choosing it too large merges
physically distinct levels (a small zero-field splitting disappears); choosing
it too small splits a symmetry-degenerate manifold into numerically distinct
pieces.  A sensible value is well below the smallest physical splitting you care
about, and well above the numerical noise of the diagonalization
($\sim10^{-12}$~eV).

\subsection{Ground-state degeneracy and excitations}

\begin{itemize}
\item \op{get\_gs\_degeneracy(tol=None)} returns the \textbf{tuple}
      $(n_{\mathrm{GS}}, E_{\mathrm{GS}})$, where $n_{\mathrm{GS}}$ is the
      integer number of states within \op{tol} of the minimum.  Take
      element \op{[0]} for the degeneracy; \op{get\_gs\_multiplicity()} returns
      that integer directly.
\item \op{get\_degeneracies()} returns the list
      $[(n_0,E_0),(n_1,E_1),\ldots]$ for all manifolds.
\item \op{get\_excitations()} returns $[E_k-E_0]$, the excitation energies of
      each manifold measured from the ground state, in eV.
\item \op{get\_manifolds()} returns the eigenvectors grouped by manifold;
      \op{get\_gs\_manifold()} returns only the ground-state group.
\end{itemize}

\subsection{Disentangling a degenerate manifold}

Inside a degenerate manifold the eigenvectors returned by the diagonalizer are
an arbitrary basis.  To get physically meaningful states one re-diagonalizes a
chosen operator $A$ within the manifold:
\begin{equation}
  A^{(\mathrm{mf})}_{ab} = \langle \psi_a | A | \psi_b\rangle ,
  \qquad \text{diagonalize } A^{(\mathrm{mf})}
  \;\Rightarrow\; \text{new basis}.
\end{equation}
\op{disentangle\_manifolds(A)} does this for every manifold and rewrites
\op{.evecs}; \op{disentangle\_gs\_manifold(A)} does it for the ground state
only.  Using $A=J_z$ (as the graphical interface does) labels the states by
their total angular-momentum projection, which is what makes the printed
wavefunctions readable.

%======================================================================
\section{Observables}
%======================================================================

\subsection{Projected expectation values}

The central tool is the representation of an operator inside a set of states,
\begin{equation}
  A^{(\mathrm{mf})}_{ab} = \langle \psi_a | A | \psi_b\rangle ,
\end{equation}
available as \op{get\_representation(A, n=6)} for the $n$ lowest states.  For
the ground-state manifold,
\begin{equation}
  \op{get\_gs\_projected\_eigenvalues}(A)
  = \mathrm{eig}\big[A^{(\mathrm{GS})}\big]
\end{equation}
returns the eigenvalues of $A$ restricted to the ground-state manifold.  This is
the quantity to look at for a degenerate ground state: for a spin doublet
projected on $S_z$ it returns $\{-\tfrac12,+\tfrac12\}$; for an orbital
projector it returns the occupations of that orbital in the manifold's natural
basis, whose mean is the average occupation.

\subsection{Angular momentum operators}

\op{Atom.Operator} contains the many-body components
\begin{equation}
  S_\alpha=\sum_i (s_\alpha)_i,\qquad
  L_\alpha=\sum_i (l_\alpha)_i,\qquad
  J_\alpha=L_\alpha+S_\alpha,
\end{equation}
under the keys \op{"sx"}, \op{"sy"}, \op{"sz"}, \op{"lx"}\ldots, \op{"jx"}\ldots,
together with the squares
\begin{equation}
  S^2=S_x^2+S_y^2+S_z^2,\qquad
  L^2=L_x^2+L_y^2+L_z^2,\qquad
  J^2=J_x^2+J_y^2+J_z^2,
\end{equation}
under \op{"s2"}, \op{"l2"}, \op{"j2"}.  From
$\langle S^2\rangle=S(S+1)$ and $\langle L^2\rangle=L(L+1)$ one reads off the
term symbol of any manifold:
\begin{equation}
  S = \tfrac{1}{2}\left(-1+\sqrt{1+4\langle S^2\rangle}\right),
  \qquad
  L = \tfrac{1}{2}\left(-1+\sqrt{1+4\langle L^2\rangle}\right).
\end{equation}

The commutation relations $[L_x,L_y]=iL_z$ etc., $[\vec L,\vec S]=0$ and
$[\vec L,\hat V_{ee}]=[\vec S,\hat V_{ee}]=0$ are verified at run time by
\op{tranci.check.check\_all(atom)}.

\subsection{Orbital projectors}

Occupations of the cubic ($t_{2g}$/$e_g$) orbitals are obtained from projectors
built out of the standard cubic harmonics,
\begin{align}
  |d_{z^2}\rangle &= |m{=}0\rangle, &
  |d_{x^2-y^2}\rangle &= \tfrac{1}{\sqrt2}\big(|{+}2\rangle+|{-}2\rangle\big), &
  |d_{xy}\rangle &= \tfrac{i}{\sqrt2}\big(|{-}2\rangle-|{+}2\rangle\big),\\
  |d_{xz}\rangle &= \tfrac{1}{\sqrt2}\big(|{-}1\rangle-|{+}1\rangle\big), &
  |d_{yz}\rangle &= \tfrac{i}{\sqrt2}\big(|{-}1\rangle+|{+}1\rangle\big). & &
\end{align}
The corresponding many-body number operators
\begin{equation}
  \hat n_{\mu}=\sum_{\sigma}\cd{\mu\sigma}\cc{\mu\sigma}
\end{equation}
are stored under the keys \op{"dz2"}, \op{"dx2y2"}, \op{"dxy"}, \op{"dxz"},
\op{"dyz"}.  Each has trace $2$ in the single-particle space (two spin
channels), and together they resolve the identity.  Additional projectors on the
magnitude of the magnetic quantum number are available as \op{"m0"}, \op{"m1"},
\op{"m2"} (projecting onto $|m|=0,1,2$), and the ten individual spin-orbital
projectors are attributes \op{atom.um2}\ldots\op{atom.dp2}.

Taking the mean of \op{get\_gs\_projected\_eigenvalues(Atom.Operator["dxy"])}
gives the average $d_{xy}$ occupation of the ground-state manifold; summing over
the five orbitals returns $n_e$.

\subsection{The $g$-tensor of a Kramers doublet}
\label{sec:gtensor}

When the ground state is a Kramers doublet --- two states, as required by
time-reversal symmetry for an odd number of electrons --- its response to a
magnetic field is that of a pseudo-spin $\tilde{S}=\tfrac12$,
\begin{equation}
  \mathcal{H}_{\mathrm{eff}}
  = \mu_B\,\vec{B}\cdot \mathbf{g}\cdot \hat{\tilde{\vec{S}}} .
\end{equation}
Projecting the magnetic moment $M_\alpha=(L+2S)_\alpha$ on the doublet gives
$2\times2$ matrices which can always be written as
$A_\alpha=\sum_\beta g_{\alpha\beta}\sigma_\beta/2$.  The individual
$g_{\alpha\beta}$ depend on the arbitrary basis chosen inside the doublet, but
the combination
\begin{equation}
  G_{\alpha\beta}=\sum_\gamma g_{\alpha\gamma}g_{\beta\gamma}
  = 2\,\mathrm{Tr}\big[A_\alpha A_\beta\big]
\end{equation}
does not.  TranCI returns the symmetric positive square root
$\mathbf{g}=\sqrt{G}$: its eigenvalues are the principal values $g_x,g_y,g_z$
and its eigenvectors the magnetic axes.

\begin{lstlisting}
M = Atom.get_manifolds(H)
g = M.get_gtensor()               # 3 x 3 symmetric g-matrix
gs, axes = M.get_principal_g()    # principal values and magnetic axes
\end{lstlisting}

\noindent
The result is exact for the doublet, not a finite difference.  Along an
arbitrary direction $\hat n$ the predicted splitting is
$\Delta E = \sqrt{\hat n\cdot \mathbf{g}^2\cdot\hat n}\;\mu_B B$, and this agrees
with the directly computed splitting to all six digits of a finite-field check
($b=10^{-6}$~eV), including for tilted (non-Cartesian) magnetic axes where the
off-diagonal elements of $G$ are of the same size as the diagonal ones.  A
purely spin doublet --- a crystal-field split $d^1$ with no spin--orbit coupling
--- gives $g_x=g_y=g_z=2$ exactly, while a low-symmetry $d^7$ doublet, e.g.
$4\hat V_{ee}+0.6\,\op{z2}+0.15\,(\op{x2}-\op{y2})+0.05\,\op{ls}$, gives strongly
anisotropic values $(0.81,\,1.70,\,9.45)$.  A ground state that is not two-fold
degenerate raises an error, since a $g$-tensor is not defined for it; the
degeneracy is judged with \op{hamiltonians.tol}, which can be overridden through
the \op{tol} keyword.

\subsection{Correlation entropy}

A single Slater determinant is fully characterized by its one-body density
matrix, whose eigenvalues (natural occupations) are then all $0$ or $1$.  Any
deviation signals genuine many-body correlation.  TranCI quantifies this with
the correlation entropy of the one-body density matrix
\begin{equation}
  \rho_{ij}=\langle \Psi |\, \cd{i}\cc{j}\, |\Psi\rangle ,
  \qquad
  \rho\,|\phi_k\rangle = n_k\,|\phi_k\rangle ,
  \qquad
  S = -\sum_k n_k \ln n_k ,
\end{equation}
computed by \op{get\_correlation\_entropy(wf)} for any many-body wavefunction
\op{wf} (typically \op{M.get\_gs\_manifold()[0]}).  $S=0$ for a single
determinant, and $S$ grows as the state becomes a superposition of many
configurations, which is exactly what a strong Coulomb interaction produces.
Occupations below $10^{-7}$ are discarded to avoid $0\ln 0$.

\subsection{Dynamical correlators}

Spectroscopies --- inelastic tunnelling, neutron scattering, electron spin
resonance --- measure the frequency-resolved response of the ion to an operator
$B$, detected through an operator $A$.  TranCI computes
\begin{equation}
  S_{AB}(\omega)
  = -\frac{1}{\pi}\,\mathrm{Im}\;
    \big\langle \Psi_0 \big| A\,
    \big[\omega + E_0 + i\delta - \mathcal{H}\big]^{-1} B
    \big|\Psi_0\big\rangle
  \;=\;\sum_n \langle \Psi_0|A|n\rangle\langle n|B|\Psi_0\rangle\,
        \mathcal{L}_\delta\!\big(\omega - (E_n-E_0)\big),
  \label{eq:dyn}
\end{equation}
where
$\mathcal{L}_\delta(x)=\frac{1}{\pi}\frac{\delta}{x^2+\delta^2}$
is a normalized Lorentzian of half-width $\delta$.  The frequency $\omega$ is
measured from the ground state, so peaks appear at the excitation energies
returned by \op{get\_excitations()}, with weights given by the matrix elements.

\begin{lstlisting}
import numpy as np
A = Atom.one2many(Atom.SP_Operator["sy"] @ Atom.SP_Operator["dz2"])
es, ds = M.get_dynamical_correlator(A, es=np.linspace(-1., 1., 100))
\end{lstlisting}

\noindent
Practical notes:
\begin{itemize}
\item \op{B} defaults to \op{A}, giving the spectral function of a single
      operator.
\item \op{delta} (default $0.03$~eV) is the Lorentzian broadening.  It must be
      larger than the spacing of the frequency grid \op{es}, or peaks will fall
      between grid points.
\item The result is \emph{averaged} over the degenerate ground states, which is
      the right thing for an unpolarized, zero-temperature measurement.
\item The default solver (\op{mode="cv"}) is the correction-vector method: for
      each frequency a shifted linear system
      $\big[(\mathcal{H}-\omega)^2+\delta^2\big]x = -\delta B|\Psi_0\rangle$
      is solved by conjugate gradients.  \op{mode="full"} inverts the resolvent
      explicitly and is a useful cross-check.  A Chebyshev (kernel polynomial)
      implementation is also provided in \op{dynamicstk}.
\end{itemize}

\subsection{Effective low-energy Hamiltonians}
\label{sec:heff}

Deep in the multiplet regime the low-energy physics of the ion is described by a
handful of states --- a pseudo-spin $S$ with $2S+1$ levels --- and by an
effective spin Hamiltonian containing anisotropy and Zeeman terms.  TranCI
extracts that Hamiltonian numerically.

The procedure is:
\begin{enumerate}
\item Project the full Hamiltonian on the $n$ lowest eigenstates,
      $h_{ab}=\langle\psi_a|\mathcal{H}|\psi_b\rangle$, and remove the trace,
      $h \rightarrow h - \frac{\mathrm{Tr}\,h}{n}\mathbb{1}$, since a constant
      shift carries no physics.
\item Project the chosen operators (by default $S_\alpha$; optionally
      $L_\alpha$, $J_\alpha$) on the same manifold, and rescale each so that its
      largest eigenvalue in magnitude equals $S=(n-1)/2$.  The projected
      operators then have the spectrum of genuine spin-$S$ components, and the
      symbols $\hat S_\alpha$ in the output mean what they say.
\item Build an operator basis from those matrices: the identity, the linear
      terms $\hat S_\alpha$, their powers $\hat S_\alpha^p$ up to $p=2$, and
      symmetrized bilinears
      $\tfrac12(\hat S_\alpha \hat S_\beta + \hat S_\beta \hat S_\alpha)$.
      Linearly dependent members are discarded.
\item Fit $h$ by least squares,
      \begin{equation}
        \min_{\{c_a\}}\;
        \Big\langle \big| h - \textstyle\sum_a c_a\,O_a \big|^2 \Big\rangle ,
      \end{equation}
      with real coefficients, using analytic gradients (\op{jax}) and $40$
      random restarts to avoid local minima.
\item Emit the result as \LaTeX, normalized to the largest coefficient so that
      the overall energy scale is factored out:
      $\mathcal{H}_{\mathrm{eff}} = c_{\max}\big[\,\hat O_1 + r_2 \hat O_2 +
      \ldots\big]$.
\end{enumerate}

The resulting Hermitian model is of the familiar form
\begin{equation}
  \mathcal{H}_{\mathrm{eff}}
  = D\,\hat S_z^2 + E\,(\hat S_x^2-\hat S_y^2)
  + \vec{g}\mu_B\vec{B}\cdot\hat{\vec{S}} + \ldots ,
\end{equation}
i.e.\ the axial and rhombic magnetic anisotropies of the ion, derived from first
principles rather than assumed.

\begin{lstlisting}
from tranci import effectivehamiltonian
deg  = M0.get_gs_multiplicity()             # size of the low-energy block
text = effectivehamiltonian.effective_spin_hamiltonian(
            M, H=H, n=deg, operators=["sx","sy","sz"])
print(text)                                  # LaTeX source
\end{lstlisting}

\noindent
\textbf{Caveats.}  The fit only means something if the chosen operator basis can
actually represent the block.  The code reports the relative error of the fitted
spectrum, prints a warning when it exceeds $1\%$, and inserts a warning in the
\LaTeX\ output.  If you see one, either enlarge the operator set (add
\op{"lx","ly","lz"}), or reconsider $n$: it should equal the multiplicity of a
manifold that is well separated from the rest of the spectrum, not an arbitrary
cut through a degenerate group.  A second entry point,
\op{effective\_hamiltonian(lowest, n, nt)}, tries the $LS$, $SJ$ and $LJ$
operator sets in turn and additionally prints the effective algebra; this is the
one used by the graphical interface.

%======================================================================
\section{Custom crystal fields and custom operators}
\label{sec:custom}
%======================================================================

The built-in Cartesian generators of Sec.~\ref{sec:cf} cover the usual point
symmetries, but any one-body operator can be supplied directly.  The route is:
write a $10\times 10$ matrix in the spin-orbital basis of Sec.~\ref{sec:spbasis},
then promote
it to the many-body space with
\begin{equation}
  \hat O = \sum_{ij} O_{ij}\, \cd{i}\cc{j}
\end{equation}
which is exactly what \op{Atom.one2many(m)} computes (with the correct fermionic
signs, via \op{edtk.states.one2many\_basis}).

To build such matrices, \op{Atom.SP\_Operator} holds the same operator
dictionary in $10\times 10$ single-particle form.  It is simply the $n_e=1$
operator set, so \op{SP\_Operator["lx"]} is the $10\times 10$ matrix $l_x$,
\op{SP\_Operator["dz2"]} the single-particle $d_{z^2}$ projector, and so on.

\begin{lstlisting}
Atom  = get_atom(ne=2)
SP    = Atom.SP_Operator

# a one-body crystal field written in single-particle language
h1 = 0.1*(SP["x4"] + SP["y4"] + SP["z4"]) + 0.3*SP["z2"] + 0.05*SP["ls"]

H = Atom.one2many(h1) + 4*Atom.Operator["Coulomb"]
M = Atom.get_manifolds(H)
\end{lstlisting}

\noindent
This is also the way to compare a one-body (single-particle) description with
the correlated many-body result: build the \emph{same} $h_1$ for \op{ne=1} and
for \op{ne=2}, and compare the orbital occupations.  The example
\op{examples/single\_occ\_VS\_many\_occ} does precisely that, and shows how the
Coulomb interaction redistributes the occupations relative to the
non-interacting filling of the same levels.

Any Hermitian $10\times 10$ matrix is legitimate: a ligand-field matrix obtained
from a density-functional calculation, a Wannier-projected on-site Hamiltonian,
or a hopping matrix.  The only requirement is that the row/column ordering
matches Sec.~\ref{sec:spbasis}.

%======================================================================
\section{Python API reference}
%======================================================================

\subsection{Entry point}

\begin{center}
\begin{tabular}{ll}
\toprule
\op{get\_atom(ne)}             & load the operator library for $n_e=1\ldots9$\\
\op{Atom.Operator}             & dict of many-body matrices ($\dim\times\dim$)\\
\op{Atom.SP\_Operator}         & dict of single-particle matrices ($10\times10$)\\
\op{Atom.one2many(m)}          & promote a $10\times10$ matrix to many body\\
\op{Atom.basis}                & list of occupation-number basis states\\
\op{Atom.get\_manifolds(H)}    & diagonalize, returns \op{Lowest\_States}\\
\op{Atom.get\_latex\_wavefunction(v)} & \LaTeX\ expansion of a wavefunction\\
\op{Atom.wavefunction\_z\_axis}& quantization axis used when printing states\\
\bottomrule
\end{tabular}
\end{center}

\subsection{Operator keys}

\begin{center}
\begin{tabular}{lll}
\toprule
Key & Operator & Notes\\
\midrule
\op{"sx" "sy" "sz"} & $S_\alpha=\sum_i (s_\alpha)_i$ & total spin\\
\op{"lx" "ly" "lz"} & $L_\alpha=\sum_i (l_\alpha)_i$ & total orbital momentum\\
\op{"jx" "jy" "jz"} & $J_\alpha=L_\alpha+S_\alpha$ & total angular momentum\\
\op{"s2" "l2" "j2"} & $S^2$, $L^2$, $J^2$ & true many-body squares\\
\op{"Coulomb"}, \op{"vc"} & $\hat V_{ee}$ & Slater--Condon, see Sec.~\ref{sec:coulomb}\\
\op{"ls"}           & $\sum_i \vec l_i\cdot\vec s_i$ & spin--orbit coupling\\
\op{"x2" "y2" "z2"} & $\sum_i (l_\alpha^2)_i$ & \emph{not} $L_\alpha^2$\\
\op{"x4" "y4" "z4"} & $\sum_i (l_\alpha^4)_i$ & cubic invariant if summed\\
\op{"x2y2"}         & $\sum_i[(l_xl_y)^2+(l_yl_x)^2]_i$ & \\
\op{"dz2" "dx2y2"} & $\hat n_\mu$ & $e_g$ occupations\\
\op{"dxy" "dxz" "dyz"} & $\hat n_\mu$ & $t_{2g}$ occupations\\
\op{"m0" "m1" "m2"} & projectors on $|m|=0,1,2$ & \\
\bottomrule
\end{tabular}
\end{center}

\subsection{\texttt{Lowest\_States}}

\begin{center}
\begin{tabular}{ll}
\toprule
\op{.evals}, \op{.evecs}, \op{.e0}, \op{.h}, \op{.atom} & spectrum and inputs\\
\op{get\_gs\_degeneracy(tol=None)}   & tuple $(n_{\mathrm{GS}},E_{\mathrm{GS}})$\\
\op{get\_gs\_multiplicity()}         & the integer degeneracy\\
\op{get\_degeneracies()}             & $[(n_k,E_k)]$ for all manifolds\\
\op{get\_excitations()}              & $[E_k-E_0]$ in eV\\
\op{get\_manifolds()}                & eigenvectors grouped by manifold\\
\op{get\_gs\_manifold()}             & ground-state eigenvectors\\
\op{get\_representation(A,n=6)}      & $\langle\psi_a|A|\psi_b\rangle$\\
\op{get\_gs\_projected\_eigenvalues(A)} & eigenvalues of $A$ in the GS manifold\\
\op{disentangle\_manifolds(A)}       & re-label all manifolds by $A$\\
\op{disentangle\_gs\_manifold(A)}    & re-label the GS manifold by $A$\\
\op{get\_correlation\_entropy(wf)}   & $-\sum_k n_k\ln n_k$\\
\op{get\_gtensor(tol=None)}          & $3\times3$ $g$-matrix of a Kramers doublet\\
\op{get\_principal\_g(tol=None)}     & principal $g$ values and magnetic axes\\
\op{get\_dynamical\_correlator(A,B=None,...)} & $(\omega, S_{AB}(\omega))$\\
\bottomrule
\end{tabular}
\end{center}

\subsection{Worked examples shipped with the code}

Each directory under \op{examples/} contains a self-contained \op{main.py},
to be run from inside that directory.

\begin{description}
\item[\op{octahedral}] Builds a cubic field directly from powers of the total
      orbital operators, $L_x^4+L_y^4+L_z^4$, adds a small uniaxial term and the
      Coulomb interaction for $n_e=1$, and prints the ground-state degeneracy
      together with the occupation of each cubic orbital.  The fastest sanity
      check of an installation (seconds).
\item[\op{orbital\_projection}] Sweeps a uniaxial crystal field for $n_e=3$ and
      plots, as a function of its strength, the five orbital occupations, the
      $S_z$ content of the ground-state manifold, and the ground-state
      degeneracy.  The total spin stays $S=3/2$ across the sweep while the
      ground-state degeneracy drops from $8$ to $4$: the crystal field quenches
      the orbital degree of freedom, not the spin.
\item[\op{single\_occ\_VS\_many\_occ}] Compares the orbital occupations of a
      correlated $d^2$ ion with those of a single electron subject to the same
      one-body field, illustrating the effect of the Coulomb interaction.  Also
      the reference for building custom crystal fields through
      \op{one2many}.
\item[\op{correlation\_entropy}] Computes the correlation entropy of the ground
      state of a $d^3$ ion with crystal field, Zeeman and spin--orbit terms.
\item[\op{dynamical\_spin\_correlator}] Computes the spin excitation spectrum of
      a $d^6$ ion, using an operator that combines a spin flip with an orbital
      projector, and plots $S_{AA}(\omega)$.
\item[\op{effective\_hamiltonian}] Fits the low-energy block of a $d^2$ ion with
      spin--orbit coupling to an effective spin model and prints the resulting
      \LaTeX\ formula.  Exercises the \op{jax}-based fitting path.
\end{description}

\noindent
The notebooks in \op{notebooks/} cover the same ground interactively:
\op{Hund\_rules.ipynb}, \op{orbital\_CF.ipynb},
\op{anisotropy\_VS\_crystal\_field.ipynb}, \op{custom\_crystal\_field.ipynb},
\op{Zeeman\_splitting.ipynb} and \op{Effective\_Hamiltonian.ipynb}.

%======================================================================
\section{The graphical interface}
%======================================================================

The graphical front-end is started with \op{tranci} on Linux and Mac
(\op{bin\textbackslash tranci.bat} on Windows) after running
\op{python install.py}, which only adds \op{bin/} to the \op{PATH}.  It needs
\op{PyQt5} and a \op{pdflatex} installation (MiKTeX on Windows, TeX~Live or
MacTeX elsewhere).

\subsection{Parameters}

All parameters are visible at once in the \op{Parameters} panel, grouped as
\op{Atom}, \op{Crystal field} and \op{Zeeman and exchange fields}.  The
interface builds a fixed Hamiltonian from the fields below.  Each field
multiplies exactly one operator of Sec.~\ref{sec:ham}; all of them are in eV
except \op{U}, which is a dimensionless multiplier (Sec.~\ref{sec:coulomb}).
Hovering over any field shows a tooltip with its meaning.  The number of
electrons and the other integer fields are spin boxes and only accept values in
their valid range; the real-valued fields reject anything that is not a number
as it is typed.

Next to every field the interface shows the operator it multiplies, written in
second quantization and rendered from \TeX\ at start-up (hovering over a
formula shows its \TeX\ source).  Here $\cd{m\sigma}$ creates an electron with
$l_z=m$ ($m=-2,\dots,2$) and spin $\sigma$ in the $d$ shell, lower-case $l,s$
are single-particle operators and $V_{ijkl}$ is the tensor of
Eq.~\eqref{eq:vee}:
\begin{center}
\begin{tabular}{ll}
\toprule
Field & Second-quantized operator\\
\midrule
\op{\# of electrons} & $\hat N=\sum_{m\sigma}\cd{m\sigma}\cc{m\sigma}$\\[2pt]
\op{U} & $U\sum_{ijkl}\sum_{\sigma\sigma'}V_{ijkl}\,\cd{i\sigma}\cd{j\sigma'}\cc{k\sigma'}\cc{l\sigma}$\\[2pt]
\op{SOC} & $\lambda\sum_{mm'\sigma\sigma'}\langle m\sigma|\vec l\cdot\vec s|m'\sigma'\rangle\,\cd{m\sigma}\cc{m'\sigma'}$\\[2pt]
\op{Uniaxial z\^{}2} & $D\sum_{m\sigma}m^{2}\,\cd{m\sigma}\cc{m\sigma}$\\[2pt]
\op{Shear x\^{}2-y\^{}2} & $E\sum_{mm'\sigma}\langle m|l_x^{2}-l_y^{2}|m'\rangle\,\cd{m\sigma}\cc{m'\sigma}$\\[2pt]
\op{Octahedral} & $O\sum_{mm'\sigma}\langle m|l_x^{4}+l_y^{4}+l_z^{4}|m'\rangle\,\cd{m\sigma}\cc{m'\sigma}$\\[2pt]
\op{Trigonal} & $t\sum_{mm'\sigma}\langle m|(\vec l\cdot\hat n)^{2}|m'\rangle\,\cd{m\sigma}\cc{m'\sigma}$, \quad $\hat n=(1,1,1)/\sqrt3$\\[2pt]
\op{Uniaxial' z\^{}4} & $D'\sum_{m\sigma}m^{4}\,\cd{m\sigma}\cc{m\sigma}$\\[2pt]
\op{Shear' (xy)\^{}2+(yx)\^{}2} & $E'\sum_{mm'\sigma}\langle m|(l_xl_y)^{2}+(l_yl_x)^{2}|m'\rangle\,\cd{m\sigma}\cc{m'\sigma}$\\[2pt]
\op{B} & $\vec B\cdot\sum_{mm'\sigma\sigma'}\langle m\sigma|\vec l+2\vec s|m'\sigma'\rangle\,\cd{m\sigma}\cc{m'\sigma'}$\\[2pt]
\op{J} & $\vec J\cdot\sum_{m\sigma\sigma'}\vec s_{\sigma\sigma'}\,\cd{m\sigma}\cc{m\sigma'}$, \quad $\vec s=\vec\sigma/2$\\[2pt]
\op{Theta}, \op{Phi} & $\hat B=(\sin\pi\theta_B\cos\pi\phi_B,\ \sin\pi\theta_B\sin\pi\phi_B,\ \cos\pi\theta_B)$, and likewise $\hat J$\\
\bottomrule
\end{tabular}
\end{center}
The uniaxial fields are diagonal in $m$ because $l_z|m\rangle=m|m\rangle$; the
matrix elements of the other one-body operators are the $10\times10$ matrices
of \op{Atom.SP\_Operator} (Sec.~\ref{sec:custom}).

\begin{center}
\begin{tabular}{lll}
\toprule
Field & Operator & Meaning\\
\midrule
\op{\# of electrons} & --- & $n_e$, from 1 to 9\\
\op{U [multiplier]}  & \op{vc} & multiplier of the Coulomb tensor (dimensionless)\\
\op{SOC [eV]}        & \op{ls} & $\lambda$, spin--orbit coupling\\
\op{Uniaxial z\^{}2} & \op{z2} & $D$, tetragonal field\\
\op{Shear x\^{}2-y\^{}2} & \op{x2-y2} & $E$, rhombic field\\
\op{Octahedral x\^{}4+y\^{}4+z\^{}4} & \op{x4+y4+z4} & $O$, cubic field ($10Dq=6O$)\\
\op{Uniaxial' z\^{}4} & \op{z4} & higher-order axial field\\
\op{Shear' (xy)\^{}2+(yx)\^{}2} & \op{x2y2} & higher-order in-plane field\\
\op{Trigonal (x+y+z)\^{}2} & rotated \op{z2} & field along $(1,1,1)$\\
\op{B [eV]}, \op{Theta\_B}, \op{Phi\_B} & $\vec b\cdot(\vec L+2\vec S)$ & Zeeman field\\
\op{J [eV]}, \op{Theta\_J}, \op{Phi\_J} & $\vec J\cdot\vec S$ & exchange field (spin only)\\
\bottomrule
\end{tabular}
\end{center}

\noindent
The angles $\theta$ and $\phi$ are given \emph{in units of $\pi$}, so
$\theta=0.5,\ \phi=0$ means a field along $x$.  The direction is
$\vec b = b(\sin\pi\theta\cos\pi\phi,\ \sin\pi\theta\sin\pi\phi,\ \cos\pi\theta)$.

The \op{Calculation} panel has three tabs.  \op{Run} holds the single-shot
calculation; its checkbox \op{Fit effective Hamiltonian with $n$ states} adds
the fit of Sec.~\ref{sec:heff} to the summary (off by default, since it
requires \op{jax}), with $n$ the dimension of the local effective model, for
example $2S+1$.  \op{Sweep} holds the parameter sweeps, with \op{Levels to
plot} limiting how many eigenvalues are drawn (0 draws all).  Under
\op{Settings}, \op{Energy tolerance [eV]} sets the degeneracy-grouping window
\op{tol} discussed above, and the \op{z axis for the wavefunctions} rotates the
quantization axis used when the wavefunctions are printed, so that a state
polarized along an arbitrary direction collapses onto a single basis ket.

The status bar at the bottom of the window reports what the interface is doing
(the progress of a sweep, whether the PDF was created, where the data were
saved); errors open a dialog.  The window is blocked while a calculation runs.

\subsection{Saving and loading parameters}

\op{File $\to$ Save parameters} (Ctrl+S) writes every field of the interface to
a JSON file, and \op{File $\to$ Load parameters} (Ctrl+O) fills the interface
back from one, so that a calculation can be reproduced or shared.  Entries of
the file that do not match a field of the current version are reported and
ignored.  Every \op{Initialize and run} also writes the parameters it used to
\op{parameters.json} next to the results, and \op{Save data} copies that file
along with them.  \op{Help $\to$ Manual} (F1) opens this document.

\subsection{Tasks}

\begin{description}
\item[Initialize and run] performs one calculation and writes
      \op{spectrum\_ci.tex}, compiling it twice with \op{pdflatex} (so that the
      table of contents is correct) into \op{spectrum\_ci.pdf}.  The summary
      contains the Hamiltonian together with a table of every parameter of
      the run, the table of manifolds with degeneracies and excitation
      energies, a table of energies with $\langle L_\alpha\rangle$ and
      $\langle S_\alpha\rangle$ for each state, the projection of all
      standard operators on the ground-state manifold (a table of expectation
      values when the ground state is non-degenerate, matrices otherwise), the
      explicit wavefunctions of the ground state and the first ten excited
      manifolds, and (optionally) the effective Hamiltonian.  The remaining
      manifolds are listed in an appendix; from Python the number of
      manifolds in the main text is the \op{nbody} argument of
      \op{write.write\_all}.
\item[Show pdf] opens that summary with the system viewer.
\item[Plot spectrum / Plot excitations] sweep the parameter selected in
      \op{Parameter to sweep} over \op{[Initial, Final]} in \op{Steps} points
      and plot all eigenvalues, referred to their centre of gravity or to the
      ground state respectively.  The entries of the sweep list carry the same
      names as the fields of the \op{Parameters} panel; the value typed in the
      swept field is ignored, all the others are kept fixed.  The angles are
      swept in units of $\pi$.
\item[Plot degeneracy] plots the ground-state degeneracy along the sweep --- the
      quickest way to locate level crossings and spin-state transitions.
\item[Plot operator] plots the eigenvalues of the operator selected next to
      the button, projected on the ground-state manifold along the sweep.
\item[Save data] copies the generated \op{*.tex}, \op{*.pdf}, \op{*.OUT} and
      \op{parameters.json} files into \op{./tranci\_data} relative to the
      directory from which \op{tranci} was launched, and reports the folder.
      It is also available as \op{File $\to$ Save data} (Ctrl+D).
\end{description}

\noindent
The interface runs inside a fresh temporary directory, so all output files are
created there until \op{Save data} is pressed.

\subsection{Output files}

\begin{center}
\begin{tabular}{ll}
\toprule
\op{spectrum\_ci.tex} / \op{.pdf} & full \LaTeX{} summary of one calculation\\
\op{parameters.json}    & every field of the interface for that calculation\\
\op{SPECTRUM.OUT}      & energy (GHz, meV) and $\langle L_\alpha\rangle$, $\langle S_\alpha\rangle$\\
\op{POPULATION.OUT}    & energy and the ten spin-orbital occupations\\
\op{EIGENVALUES.OUT}   & swept parameter vs.\ eigenvalues\\
\op{DEGENERACY.OUT}    & swept parameter vs.\ ground-state degeneracy\\
\op{OPERATOR\_VALUES.OUT} & swept parameter vs.\ projected operator eigenvalues\\
\bottomrule
\end{tabular}
\end{center}

%======================================================================
\section{Validation}
\label{sec:validation}
%======================================================================

\subsection{Hund's first and second rules}

A direct test of the Coulomb tensor and of the many-body machinery is that the
interaction alone, with no crystal field and no spin--orbit coupling, must
reproduce the free-ion ground terms predicted by the first two Hund rules:
maximize $S$, then maximize $L$.  Diagonalizing $\hat V_{ee}$ by itself and
reading off $\langle S^2\rangle$ and $\langle L^2\rangle$ in the ground-state
manifold gives

\begin{center}
\begin{tabular}{ccccccc}
\toprule
$n_e$ & Hilbert dim. & $S$ & $L$ & term & degeneracy & $(2S{+}1)(2L{+}1)$\\
\midrule
1 & 10  & 1/2 & 2 & $^2D$ & 10 & 10\\
2 & 45  & 1   & 3 & $^3F$ & 21 & 21\\
3 & 120 & 3/2 & 3 & $^4F$ & 28 & 28\\
4 & 210 & 2   & 2 & $^5D$ & 25 & 25\\
5 & 252 & 5/2 & 0 & $^6S$ & 6  & 6\\
6 & 210 & 2   & 2 & $^5D$ & 25 & 25\\
7 & 120 & 3/2 & 3 & $^4F$ & 28 & 28\\
8 & 45  & 1   & 3 & $^3F$ & 21 & 21\\
9 & 10  & 1/2 & 2 & $^2D$ & 10 & 10\\
\bottomrule
\end{tabular}
\end{center}

\noindent
Every entry matches the textbook free-ion term, including the $L=0$ half-filled
$^6S$ case and the particle--hole symmetry between $n_e$ and $10-n_e$.  The
degeneracies come out exactly as $(2S+1)(2L+1)$, confirming that the ground
manifold is a single $LS$ term.  This calculation is reproduced by

\begin{lstlisting}
from tranci import hamiltonians
hamiltonians.tol = 1e-6
Atom = get_atom(ne=2)
M = Atom.get_manifolds(Atom.Operator["Coulomb"])
gs = M.get_gs_manifold()
s2 = hamiltonians.get_projected_eigenvalues(gs, Atom.Operator["s2"]).mean()
l2 = hamiltonians.get_projected_eigenvalues(gs, Atom.Operator["l2"]).mean()
\end{lstlisting}

\subsection{The $d^2$ multiplet spectrum}

A sharper test uses the full excited-term structure.  For $n_e=2$ the Coulomb
interaction alone must split the $45$ states into the five free-ion terms
$^3F$, $^1D$, $^3P$, $^1G$, $^1S$, at energies given exactly by the Racah
parameters.  With the multiplier set to $u=1$, \op{get\_excitations()} returns

\begin{center}
\begin{tabular}{lccc}
\toprule
Term & Degeneracy & Racah expression & $\Delta E$ [eV]\\
\midrule
$^3F$ & 21 & $0$          & $0$\\
$^1D$ & 5  & $5B+2C$      & $0.68488$\\
$^3P$ & 9  & $15B$        & $0.79301$\\
$^1G$ & 9  & $12B+2C$     & $1.05495$\\
$^1S$ & 1  & $22B+7C$     & $2.63497$\\
\bottomrule
\end{tabular}
\end{center}

\noindent
The four excitation energies are reproduced to every printed digit by the single
pair $B=52.87$~meV, $C=210.3$~meV quoted in Sec.~\ref{sec:coulomb}, and the
degeneracies $21,5,9,9,1$ sum to $45$.  This simultaneously validates the
Coulomb tensor, the fermionic signs of the many-body construction and the
degeneracy grouping.

\subsection{Hund's third rule}

Hund's \emph{third} rule is not contained in $\hat V_{ee}$; it follows from the
spin--orbit term.  Adding $\lambda\,\op{ls}$ with $\lambda>0$ --- the physical
one-electron constant, positive for every filling --- yields $J=|L-S|$ below
half filling and $J=L+S$ above it, as tabulated in Sec.~\ref{sec:soc}.  No sign
change of the input is needed.

\subsection{Algebraic consistency}

Additional internal consistency checks are performed by
\op{tranci.check.check\_all(atom)}, which verifies the angular-momentum algebra
and the commutation of $\vec L$ and $\vec S$ with the Coulomb operator.

%======================================================================
\section{Practical notes and limitations}
%======================================================================

\begin{itemize}
\item \textbf{Only the $d$ shell.}  The model contains no ligand orbitals and no
      charge fluctuations: the electron number $n_e$ is fixed.  Charge-transfer
      physics, ligand-to-metal hybridization and covalency are represented only
      indirectly, through the effective crystal-field parameters and through a
      reduced (nephelauxetic) Coulomb multiplier.

\item \textbf{The Coulomb prefactor is dimensionless.}  See
      Sec.~\ref{sec:coulomb}.  Only the ratios $B$ and $C$ affect the multiplet
      splittings at fixed $n_e$.

\item \textbf{Crystal-field keys are sums of one-body operators.}
      $\op{x2}\neq L_x^2$.  Use \op{"l2"}, \op{"s2"}, \op{"j2"} when you want
      the many-body squares.

\item \textbf{Degeneracy is a numerical statement.}  It depends on
      \op{hamiltonians.tol}.  Always set it deliberately when a small splitting
      matters.

\item \textbf{Dense diagonalization.}  Every call to \op{get\_manifolds} builds a
      dense matrix and computes the full spectrum, at $O(\dim^3)$ cost.  This is
      negligible for a single point, but a sweep with many steps repeats it at
      every step; reuse the \op{Atom} object (loading it re-reads files from
      disk) and keep sweeps modest.

\item \textbf{Silent zeros on malformed files.}  \op{read.read\_matrix} wraps its
      parsing in a bare \op{except} and returns a zero matrix if anything goes
      wrong, rather than raising.  A corrupted or truncated \op{.op} file
      therefore shows up as a missing term in the Hamiltonian, not as an error.
      If a spectrum looks unexpectedly degenerate, verify the operator library.

\item \textbf{The $g$-tensor needs a Kramers doublet.}  \op{get\_gtensor} is
      defined only for a two-fold degenerate ground state and raises otherwise
      (Sec.~\ref{sec:gtensor}).  For a larger manifold, read the anisotropy off
      the effective Hamiltonian of Sec.~\ref{sec:heff} instead.

\item \textbf{Effective-Hamiltonian fits need checking.}  The fit is a numerical
      least-squares problem in a chosen operator basis; always look at the
      reported relative error before quoting the coefficients.

\item \textbf{Hyperfine coupling is library-only.}  It is not reachable from the
      graphical interface, and only spin-$\tfrac12$ nuclei are implemented.

\item \textbf{Dead code.}  \op{src/tranci.py} (the old GTK front-end),
      \op{src/tranci/states.py} and \op{src/tranci/fix\_coulomb.py} are legacy
      and are not imported by anything.  \op{interface\_pyqt/interface.py} is
      generated from \op{interface.ui} and must not be edited by hand.
\end{itemize}

\end{document}
